\documentclass[10pt,a4paper]{amsart}
\usepackage{amsmath,amssymb,amsthm,graphicx}
\usepackage{verbatim}
\newenvironment{reference}{\comment}{\endcomment}
%\newenvironment{reference}{}{}
\newenvironment{slogan}{\comment}{\endcomment}
\newenvironment{history}{\comment}{\endcomment}

% For commutative diagrams we use Xy-pic
\usepackage[all]{xy}

% We use 2cell for 2-commutative diagrams.
\xyoption{2cell}
% Theorem environments.
%
\theoremstyle{plain}
\newtheorem{theorem}[subsection]{Theorem}
\newtheorem{proposition}[subsection]{Proposition}
\newtheorem{lemma}[subsection]{Lemma}

\theoremstyle{definition}
\newtheorem{definition}[subsection]{Definition}
\newtheorem{example}[subsection]{Example}
\newtheorem{exercise}[subsection]{Exercise}
\newtheorem{situation}[subsection]{Situation}

\theoremstyle{remark}
\newtheorem{remark}[subsection]{Remark}
\newtheorem{remarks}[subsection]{Remarks}

\numberwithin{equation}{subsection}

% Macros
%
\def\lim{\mathop{\mathrm{lim}}\nolimits}
\def\colim{\mathop{\mathrm{colim}}\nolimits}
\def\Spec{\mathop{\mathrm{Spec}}}
\def\Hom{\mathop{\mathrm{Hom}}\nolimits}
\def\Ext{\mathop{\mathrm{Ext}}\nolimits}
\def\SheafHom{\mathop{\mathcal{H}\!\mathit{om}}\nolimits}
\def\SheafExt{\mathop{\mathcal{E}\!\mathit{xt}}\nolimits}
\def\Sch{\mathit{Sch}}
\def\Mor{\mathop{\mathrm{Mor}}\nolimits}
\def\Ob{\mathop{\mathrm{Ob}}\nolimits}
\def\Sh{\mathop{\mathit{Sh}}\nolimits}
\def\NL{\mathop{N\!L}\nolimits}
\def\CH{\mathop{\mathrm{CH}}\nolimits}
\def\proetale{{pro\text{-}\acute{e}tale}}
\def\etale{{\acute{e}tale}}
\def\QCoh{\mathit{QCoh}}
\def\Ker{\mathop{\mathrm{Ker}}}
\def\Im{\mathop{\mathrm{Im}}}
\def\Coker{\mathop{\mathrm{Coker}}}
\def\Coim{\mathop{\mathrm{Coim}}}

% Boxtimes
%
\DeclareMathSymbol{\boxtimes}{\mathbin}{AMSa}{"02}

%
% Macros for moduli stacks/spaces
%
\def\QCohstack{\mathcal{QC}\!\mathit{oh}}
\def\Cohstack{\mathcal{C}\!\mathit{oh}}
\def\Spacesstack{\mathcal{S}\!\mathit{paces}}
\def\Quotfunctor{\mathrm{Quot}}
\def\Hilbfunctor{\mathrm{Hilb}}
\def\Curvesstack{\mathcal{C}\!\mathit{urves}}
\def\Polarizedstack{\mathcal{P}\!\mathit{olarized}}
\def\Complexesstack{\mathcal{C}\!\mathit{omplexes}}
% \Pic is the operator that assigns to X its picard group, usage \Pic(X)
% \Picardstack_{X/B} denotes the Picard stack of X over B
% \Picardfunctor_{X/B} denotes the Picard functor of X over B
\def\Pic{\mathop{\mathrm{Pic}}\nolimits}
\def\Picardstack{\mathcal{P}\!\mathit{ic}}
\def\Picardfunctor{\mathrm{Pic}}
\def\Deformationcategory{\mathcal{D}\!\mathit{ef}}
\usepackage{cleveref}
\crefname{originallemma}{original lemma}{original lemmas}
\Crefname{originallemma}{Original lemma}{Original lemmas}
\makeatletter
\newlabel{lemma-out-of-finite}{{05NR}{source}{}{sourcepointer-1}{}}
\newlabel{lemma-out-of-finite@cref}{{[originallemma][0][]05NR}{[1][0][]source}{}{}{}}
\makeatother
\crefname{originalequation}{original equation}{original equations}
\Crefname{originalequation}{Original equation}{Original equations}
\makeatletter
\newlabel{equation-compare}{{05NN}{source}{}{sourcepointer-2}{}}
\newlabel{equation-compare@cref}{{[originalequation][0][]05NN}{[1][0][]source}{}{}{}}
\makeatother
\title{Modules are small with respect to injections for ordinals of large cofinality}
\author{The Stacks Project Authors}
\date{stacks-project:injectives.tex:local-snapshot-20261012; GNU FDL 1.2 or later}
\begin{document}
\maketitle
\noindent\textit{Editorial note.} All delivered mathematical text is the author's own, copied from the pinned Stacks Project chapter file sites-cohomology.tex; no statement or proof step was rewritten, reordered or removed.\par\medskip
\noindent\textit{Editorial note.} The proof is the author's complete argument; the opening remark that the proof is straightforward is followed by the argument itself, including the special case and the cofinality bound, so no step is left to the reader.\par\medskip
\noindent\textit{Editorial note.} Every labelled reference occurring inside the counted spans is carried in assembly.external\_references as a non-load-bearing navigation pointer (same-chapter ones at their definition site, cross-chapter ones at the invocation inside the counted span), so no load-bearing mathematics is replaced by a stub. No bibliography citation occurs inside the counted spans.\par\medskip
\noindent\textit{Editorial note.} The same-name check for this card scanned every staged Stacks chapter of this repository -- all sources/stacks-*/pinned-original-source.tex files, the card's own chapter excluded -- and the labels found are recorded in dedup.same\_name\_cross\_chapter on the card.\par\medskip
\noindent\textit{Editorial note.} The preamble is the chapter's own retained minimal preamble (sources/stacks-sites-cohomology/retained-preamble.tex), a verbatim excerpt of the native preamble: the Xy-pic setup and the native theorem declarations with the mathematical macros.\par\medskip
\noindent\textit{Allowed standard prerequisites (editorial).}\begin{itemize}
\item Smallness with respect to injections. Same chapter, the definition of an object being alpha-small with respect to a class of morphisms. Supplies the property the proposition establishes, namely that every map from the module into a colimit of an inductive system factors through one stage.
\item The finite case. Same chapter, Lemma lemma-out-of-finite. Supplies the special case of a finite module, which is treated first and which the general argument imitates.
\item Cofinality and upper bounds in an ordinal. Set theory, the elementary facts about cofinality of an ordinal and upper bounds of small subsets. Supplies that a subset of an ordinal of cardinality at most kappa has an upper bound strictly below an ordinal of cofinality bigger than kappa.
\end{itemize}\medskip
\ifdefined\hypertarget\hypertarget{sourcepointer-1}{}\fi
\noindent\textit{Source reference.} Original lemma 05NR: sources/stacks-injectives/pinned-original-source.tex, source lines 326--326. Navigation pointer invoked at proof line 400; the counted span contains the pointer only and the cited result is not reproduced inside it, so the pointer is registered as a non-load-bearing navigation stub in the established form for Stacks material. No bibliography citation is involved.\par\medskip
\ifdefined\hypertarget\hypertarget{sourcepointer-2}{}\fi
\noindent\textit{Source reference.} Original equation 05NN: sources/stacks-injectives/pinned-original-source.tex, source lines 279--279. Navigation pointer invoked at proof line 327; the counted span contains the pointer only and the cited result is not reproduced inside it, so the pointer is registered as a non-load-bearing navigation stub in the established form for Stacks material. No bibliography citation is involved.\par\medskip
\begin{proposition}
\label{proposition-modules-are-small}
Let $R$ be a ring. Let $M$ be an $R$-module.
Let $\kappa$ the cardinality of the set of submodules of $M$.
If $\alpha$ is an ordinal whose cofinality is bigger than $\kappa$,
then $M$ is $\alpha$-small with respect to injections.
\end{proposition}

\begin{proof}
The proof is straightforward, but let us first think about a special case.
If $M$ is finite, then the claim is that for any inductive system
$\{B_\beta\}$ with injections between them, parametrized by a
limit ordinal, any map $M \to \colim B_\beta$ factors through one of
the $B_\beta$. And this we proved in
Lemma \ref{lemma-out-of-finite}.

\medskip\noindent
Now we start the proof in the general case.
We need only show that the map (\ref{equation-compare}) is a surjection.
Let $f : M \to \colim B_\beta$ be a map.
Consider the subobjects $\{f^{-1}(B_\beta)\}$ of $M$, where $B_\beta$
is considered as a subobject of the colimit $B = \bigcup_\beta B_\beta$.
If one of these, say $f^{-1}(B_\beta)$, fills $M$,
then the map factors through $B_\beta$.

\medskip\noindent
So suppose to the contrary that all of the $f^{-1}(B_\beta)$ were proper
subobjects of $M$. However, we know that
$$
\bigcup f^{-1}(B_\beta) = f^{-1}\left(\bigcup B_\beta\right) = M.
$$
Now there are at most $\kappa$ different subobjects of $M$ that occur among
the $f^{-1}(B_\alpha)$, by hypothesis.
Thus we can find a subset $S \subset \alpha$ of cardinality at most
$\kappa$ such that as $\beta'$ ranges over $S$, the
$f^{-1}(B_{\beta'})$ range over \emph{all} the $f^{-1}(B_\alpha)$.

\medskip\noindent
However, $S$ has an upper bound $\widetilde{\alpha} < \alpha$ as
$\alpha$ has cofinality bigger than $\kappa$. In particular, all the
$f^{-1}(B_{\beta'})$, $\beta' \in S$ are contained in
$f^{-1}(B_{\widetilde{\alpha}})$.
It follows that $f^{-1}(B_{\widetilde{\alpha}}) = M$.
In particular, the map $f$ factors through $B_{\widetilde{\alpha}}$.
\end{proof}

\end{document}
